\section{第1章の解答}\label{解答:第1章}
\paragraph{1.1}
平均の速度は,位置の変化を時間の変化で割って
\begin{align*}
 \frac{s(3+h)-s(3)}h
 &=\frac{(3+h)^2-9}{h}\\
 &=\frac{6h+h^2}{h}\\
 &=6+h \qquad(h\ne0)
\end{align*}
である.\,$h=0.2$ のとき $6.2$,\,$h=-0.2$ のとき $5.8$ となる.\,$h$ を正負どちらから $0$ に近づけても $6+h$ は $6$ に近づくので,瞬間の速度は毎秒 $6$ メートルと予想できる.ここで扱ったように,\,$0$ に近い負の $h$ を取ったときも分子と分母がともに負となるため,同じ平均変化率の式を使える.

\paragraph{1.2}
台形の平行な2辺の長さは $1$ と $5$,その間の距離は $2$ だから,面積は
\[
 \frac{1+5}{2}\times2=6
\]
である.\,$n$ 等分した幅は $\frac2n$,左端の座標は $\frac{2k}{n}$（$k=0,\ldots,n-1$）なので,
\begin{align*}
 L_n
 &=\frac2n\sum_{k=0}^{n-1}\left(1+\frac{4k}{n}\right)\\
 &=\frac2n\left\{n+\frac4n\cdot\frac{n(n-1)}2\right\}\\
 &=6-\frac4n.
\end{align*}
右端については添字を $1$ から $n$ までとし,
\begin{align*}
 U_n
 &=\frac2n\sum_{k=1}^{n}\left(1+\frac{4k}{n}\right)\\
 &=\frac2n\left\{n+\frac4n\cdot\frac{n(n+1)}2\right\}\\
 &=6+\frac4n.
\end{align*}
$\frac4n$ が $0$ に近づくので,どちらも $6$ に近づく.

\paragraph{1.3}
各小区間の中央は $\frac18,\frac38,\frac58,\frac78$ である.幅はすべて $\frac14$ だから,求める和は
\begin{align*}
 M_4&=\frac14\left\{\left(\frac18\right)^2
 +\left(\frac38\right)^2+\left(\frac58\right)^2
 +\left(\frac78\right)^2\right\}\\
 &=\frac{1+9+25+49}{256}\\
 &=\frac{21}{64}.
\end{align*}
一方,
\[
 L_4=\frac14\left(0+\frac1{16}+\frac4{16}+\frac9{16}\right)=\frac7{32},
 \qquad
 U_4=\frac14\left(\frac1{16}+\frac4{16}+\frac9{16}+1\right)=\frac{15}{32}.
\]
よって $L_4=\frac{14}{64}<M_4=\frac{21}{64}<\frac{30}{64}=U_4$ である.中央の高さは各区間の左端と右端の高さの間にあるため,この大小関係は計算前にも分かる.ただし,中央の和が真の面積より大きいか小さいかは,この大小関係だけでは決まらない.

\paragraph{1.4}
割線の傾きは
\begin{align*}
 \frac{(1+h)^3-1}{h}
 &=\frac{3h+3h^2+h^3}{h}\\
 &=3+3h+h^2.
\end{align*}
最後の約分には $h\ne0$ を使った.\,$h$ を $0$ に近づけると $3h,h^2$ はともに $0$ に近づくので,接線の傾きは $3$ と予想できる.

\paragraph{1.5}
$4$ 等分した幅は $\frac12$ であるから,
\begin{align*}
 L_4&=\frac12\left\{0^2+\left(\frac12\right)^2+1^2+\left(\frac32\right)^2\right\}=\frac74,\\
 U_4&=\frac12\left\{\left(\frac12\right)^2+1^2+\left(\frac32\right)^2+2^2\right\}=\frac{15}{4}.
\end{align*}
したがって $\frac74\le S\le\frac{15}{4}$ である.一般の $n$ 等分では,幅は $\frac2n$,高さは $\left(\frac{2k}{n}\right)^2$ となる.幅が $2$ 倍,高さが $4$ 倍なので,長方形の面積もその和も $8$ 倍になる.\,$0\le x\le1$ の場合の和が $\frac13$ に近づくことから,\,$S=\frac83$ と予想できる.横幅を $2$ 倍にしたから面積も $2$ 倍,とはならない.高さの変化も含める必要がある.

\paragraph{1.6}
位置の変化は速度に符号を付けたまま時間を掛けて足すので,
\[
 3\times4+(-2)\times3=6
\]
メートルである.道のりは移動した長さを足すので,
\[
 3\times4+2\times3=18
\]
メートルである.後半の速度は負であり,前半に進んだ向きとは逆に $6$ メートル戻っている.位置の変化ではこの分を引くが,道のりではこの移動も正の長さとして足す.

\paragraph{1.7}
左端の和と右端の和は
\[
 L_n=\frac{T^3}{n^3}\sum_{k=0}^{n-1}k^2
 =T^3\left(\frac13-\frac1{2n}+\frac1{6n^2}\right),
\]
\[
 U_n=\frac{T^3}{n^3}\sum_{k=1}^{n}k^2
 =T^3\left(\frac13+\frac1{2n}+\frac1{6n^2}\right)
\]
である.\,$t\ge0$ で速度 $t^2$ は増加するので,位置の変化を $D$ とすると $L_n\le D\le U_n$ が成り立つ.どちらも$\frac{T^3}{3}$へ近づくので,位置の変化は$\frac{T^3}{3}$である.速度の単位と時間の単位を掛けて位置の単位になることも確かめておこう.
\paragraph{1.8}
有限等比数列の和から$P_N=\frac16\{1-4^{-(N+1)}\}$であり,\,$\frac16-P_N=\frac1{6\cdot4^{N+1}}$となる.\,$N=2$では$\frac1{384}>\frac1{1000}$,\,$N=3$では$\frac1{1536}<\frac1{1000}$である.差は$N$とともに減少するので,最小の$N$は$3$となる.
